\section{Downstream Analysis: \lme{} (Non-Final Evidence)}
\label{sec:downstream}

To probe whether write-time governance helps a long-horizon task, we used a
stratified $180$-item slice of \lme{} ($6$ competencies $\times\,30$, seed $0$).
No run on this track is eligible for a final comparative claim. We therefore use
the exact ledger labels --- \texttt{diagnostic\_only},
\texttt{failure\_replay\_only}, \texttt{incomplete}, and
\texttt{externally\_pending} --- and do not collapse them into a generic result.

\subsection{The contaminated v7 diagnostic}

The v7 controller run is provider-contaminated: it accumulated $316$ HTTP-402
credit-exhaustion traces, $388$ empty-extraction fallbacks, at least $12$ manual
relaunches, and a no-op escalation tier (the ``big'' and ``fast'' models were
identical). It is therefore \texttt{diagnostic\_only}. On this run, naive
overwrite (B4, $114/180=63.3\%$) outperformed the controller
($110/180=61.1\%$). This does not establish general superiority for either
system, but it rules out a claim that the available downstream evidence shows a
controller advantage.

\begin{table}[t]
\centering
\small
\caption{Non-final \lme{} evidence. Replay and top-up rows are selected subsets,
not full reruns; canonical statuses come from ledger version 1.1.}
\label{tab:lme}
\begin{tabular}{@{}l r r l@{}}
\toprule
Run or construction & correct/$n$ & accuracy & ledger status \\
\midrule
v7 B1 no-memory        & 9/180   & 0.0500 & \texttt{diagnostic\_only} \\
v7 B2 flat-vector      & 80/180  & 0.4444 & \texttt{diagnostic\_only} \\
v7 B3 recency          & 43/180  & 0.2389 & \texttt{diagnostic\_only} \\
v7 B4 naive-overwrite  & 114/180 & 0.6333 & \texttt{diagnostic\_only} \\
v7 \ask{} controller   & 110/180 & 0.6111 & \texttt{diagnostic\_only} \\
v8 replay of v7 failures & 25/70 & 0.3571 & \texttt{failure\_replay\_only} \\
v8 carry-forward construction & 139/180 & 0.7722 & \texttt{failure\_replay\_only} \\
v8 top-up of former successes & 71/84 & 0.8452 & \texttt{incomplete} \\
replay + top-up, raw observed & 96/154 & 0.6234 & \texttt{incomplete} \\
homogeneous \texttt{final-1.1} & --/180 & -- & \texttt{externally\_pending} \\
\bottomrule
\end{tabular}
\end{table}

\subsection{Why the 77.2\% construction is ineligible}

The v8 replay selected exactly the $70$ v7 controller failures and repaired
$25/70$. It was not a rerun of the full slice. The reported $139/180=77.2\%$
then combined $110$ carried-forward v7 successes ($108$ original
Gemma-judged successes plus two manual v7 adjudications), $25$ Qwen replay
repairs, and four later Codex/GPT-5 overrides. Before the four overrides, the
same construction was $135/180=75.0\%$. Judge identity and pipeline version are
therefore confounded with whether an item had previously failed, and the fixes
were evaluated on the failures that motivated them. The ledger correctly labels
both figures \texttt{failure\_replay\_only}, not final or diagnostic full-run
accuracy.

The top-up is non-empty and directly contradicts the carry-forward assumption.
It re-executed $84/110$ former v7 successes under the v8 attempt: $71$ remained
correct and $13$ regressed. Combining those $71$ with the $25$ raw replay
repairs gives $96/154$ observed. Treating the $26$ unresolved items as all wrong
or all correct gives full-slice raw bounds of $53.3$--$67.8\%$. Retaining the
four replay overrides gives bounds of $55.6$--$70.0\%$. Even the optimistic
bound is below $77.2\%$, so this attempt cannot reproduce the carry-forward
construction. These are missing-data bounds, not confidence intervals.

\subsection{Single-session preference: carry-forward artifact and taxonomy}

Single-session preference (SSP) was the weakest v7 competency at $11/30$. That
is a contaminated diagnostic run only, not a clean v8 measurement.
The replay targeted the $19$ v7 SSP failures and repaired $0/19$; the top-up
completed $10/11$ former SSP successes and retained only $8/10$. Across the two
selected subsets, $8/29$ completed SSP items are correct and one is unresolved,
so the maximum is $9/30$. The earlier statement that SSP stayed ``unchanged at
$11/30$'' is only a carry-forward artifact.

The replay and v7 prompt contexts differ, so persistence of the $19$ failures
does \emph{not} isolate the answer model or support a same-context causal claim.
It does, however, identify a concrete set for descriptive error analysis. By
inspecting whether each decisive gold fact appeared in the constructed prompt,
we classified the $19$ v7 failures as shown in Table~\ref{tab:ssp}.
Grounding-supply failures (retrieval omission plus cross-session contamination)
are primary in $10/19$ and present in $15/19$; answer-side failures are primary
in $9/19$. This taxonomy describes the inspected artifacts and does not prove
that the write policy is causally exonerated.

\begin{table}[t]
\centering
\small
\caption{Primary-category taxonomy of the $19$ v7 SSP failures, from the
released taxonomy artifact.}
\label{tab:ssp}
\begin{tabular}{@{}l r r@{}}
\toprule
Primary failure category & Count & Share \\
\midrule
retrieval-omission                         & 6 & 31.6\% \\
cross-session-contamination                & 4 & 21.1\% \\
target-attribute-extraction-miss           & 3 & 15.8\% \\
distractor-entity-binding                  & 2 & 10.5\% \\
polarity-error                             & 2 & 10.5\% \\
unsupported-answer-specificity             & 1 & 5.3\% \\
immediate-request-vs-enduring-preference   & 1 & 5.3\% \\
\midrule
\multicolumn{3}{@{}l}{\emph{Grounding-supply family: primary $10/19$, present $15/19$}} \\
\multicolumn{3}{@{}l}{\emph{Answer-side family: primary $9/19$}} \\
\bottomrule
\end{tabular}
\end{table}

\subsection{Proposed intervention and pending homogeneous run}

The taxonomy motivates a proposed but unexecuted intervention: represent each
preference as a typed record (holder/target/attribute/polarity/scope/stability/
source), require retrieval to surface the supporting-session record for the
query subject, and constrain answers to retrieved, grounded preferences while
honoring explicit negatives. We make no accuracy claim for this design.

The only route to an eligible downstream result is a homogeneous full-slice run.
The frozen final protocol (version 1.1) targets all $180$ items with one answer
model and logged judge routing. Its two-item smoke test stopped before inference:
the supplied OpenRouter credential returned HTTP 403 even on a minimal Qwen
request. The ledger therefore marks the run \texttt{externally\_pending}; no
score is estimated or carried forward.
